\seksjon{3.2}{Andregradsfunksjoner}
\undertittel{Parabelen, symmetrilinjen og toppunktet.}
\maal{Finne nullpunkter, symmetrilinje og ekstremalpunkt til andregradsfunksjoner, og bruke fullstendige kvadraters metode i problemløsing}
\kompetansemaal{6,9}

\hvakanjeg{
\kan{1}{avgjøre om parabelen vender opp eller ned, og finne skjæringen med $y$-aksen}
\kan{1}{finne nullpunkter og symmetrilinjen til $f(x) = ax^2 + bx + c$}
\kan{2}{finne bunnpunkt eller toppunkt ved å sette symmetrilinjen inn i $f$}
\kan{2}{forklare hvorfor symmetrilinjen er $x = -\frac{b}{2a}$}
\kanskille
\kan{1}{skrive $x^2 + bx + c$ på formen $(x - p)^2 + q$}
\kan{2}{skrive $ax^2 + bx + c$ på toppunktform $a(x - p)^2 + q$}
\kan{2}{lese av ekstremalpunkt og verdimengde fra toppunktformen}
\kan{2}{finne største og minste verdi i praktiske problemer}
}
